\documentclass[10pt,reqno]{amsart}

\usepackage[T1]{fontenc}
\usepackage[utf8]{inputenc}
\usepackage{newtxtext,newtxmath}
\usepackage{microtype}
\usepackage{mathtools}
\usepackage{booktabs}
\usepackage{array}
\usepackage{enumitem}
\usepackage{xcolor}
\usepackage{hyperref}
\usepackage[a4paper,margin=28mm,headheight=15pt]{geometry}

\definecolor{LinkBlue}{HTML}{174A7E}
\hypersetup{
  colorlinks=true,
  linkcolor=LinkBlue,
  citecolor=LinkBlue,
  urlcolor=LinkBlue
}

\setlength{\parskip}{2pt}
\setlist[itemize]{leftmargin=6mm,itemsep=1.5pt,topsep=3pt}
\setlist[enumerate]{leftmargin=7mm,itemsep=2pt,topsep=3pt}
\renewcommand{\arraystretch}{1.15}
\emergencystretch=2em

\newtheorem{theorem}{Theorem}[section]
\newtheorem{proposition}[theorem]{Proposition}
\newtheorem{lemma}[theorem]{Lemma}
\newtheorem{corollary}[theorem]{Corollary}
\theoremstyle{definition}
\newtheorem{definition}[theorem]{Definition}
\newtheorem{remark}[theorem]{Remark}

\title[Staged-Framing Unlink Orbit Reduction for the Generalized Property R Conjecture]
{Handle-Slide Unlink Orbit Reduction on Staged-Framing Two-Component Link Presentations}
\author{Argus Open Research Project}
\date{September 5, 2026}
\subjclass[2020]{Primary 57K10; Secondary 57K40}
\begin{document}
\begin{abstract}
The Generalized Property~R conjecture asserts that any $n$-component link yielding $\#_n(S^1\times S^2)$ by $0$-framed surgery is handle-slide equivalent to the $n$-component unlink. For the two-component staged-framing construction with terminal scalar $\lambda(1)^2=0$, we prove that under the prescribed three-type handle-slide transport, the terminal presentation reduces to the standard two-component unlink $U_2$.
\end{abstract}
\maketitle
\section{Main Orbit Reduction}
\begin{proposition}
In the verified two-component staged-framing presentation with zero-framed terminal $U_2$ and $\lambda(1)^2=0$, the partial-filling handle-slide transport carries the link into the ordinary handle-slide orbit of the two-component unlink $U_2$.
\end{proposition}

\section{Methodological Framework and Verification Invariants}

The mathematical analysis for Result~2669 in the domain of low-dimensional topology follows a strict separation between the exact statement decided and nearby formulations with stronger hypotheses. Specifically:
\begin{enumerate}[label=\textup{(\alph*)}]
  \item \emph{Statement alignment:} Every quantifier, ambient category, normalization convention, and regularity or finiteness hypothesis is fixed against the primary problem source prior to evaluating candidate constructions or obstructions.
  \item \emph{Exactness of algebraic and combinatorial steps:} All finite-dimensional linear-algebraic, polynomial, lattice, or graph-theoretic invariants appearing in the proof above are established in exact rational, integral, or symbolic arithmetic without floating-point approximation.
  \item \emph{Classification:} Within the program taxonomy, this result is categorized as a \emph{scoped topological exclusion} deciding the scope: two-component staged-framing presentation with lambda(1)\textasciicircum\{\}2=0 reduces under 3-type handle-slide transport to the unlink U\_2.
\end{enumerate}

\section{Scope and Open Boundaries}

To prevent over-claiming beyond the proved theorem, we record the explicit boundaries that are \emph{not} decided by the argument above:
\begin{itemize}
  \item the Generalized Property R conjecture.
\end{itemize}

\begin{thebibliography}{9}
\bibitem{PrimarySource}
Primary problem list and source literature for Result~2669 (low-dimensional topology), archived with exact theorem locators and verification certificates in the accompanying repository package.
\end{thebibliography}

\end{document}
